\documentclass[12pt,a4paper]{article}
\usepackage[utf8]{inputenc}
\usepackage[spanish]{babel}
\usepackage{geometry}
\geometry{top=2.5cm, bottom=2.5cm, left=3cm, right=3cm}
\usepackage{pdfpages}
\usepackage{hyperref}
\usepackage{url}


\title{Evolución Histórica de la Inteligencia Artificial}
\author{Diego Ramírez Ibarra}
\date{Septiembre 2026}

\begin{document}

\maketitle

\section{Introduccion}
La historia de la Inteligencia Artificial (IA) es, en esencia, la crónica de una asimetría fundamental: la brecha entre la imaginación teórica y la realidad material. Durante gran parte del siglo XX, las matemáticas y la teoría computacional que fundamentan las redes neuronales artificiales avanzaron a un ritmo acelerado, estableciendo paradigmas que prometían replicar el razonamiento humano. Sin embargo, estas arquitecturas permanecieron latentes, atrapadas en un letargo impuesto por las severas limitaciones del hardware de la época. Alan Turing, en sus trabajos fundamentales, ya vislumbraba el potencial de las máquinas para aprender de la experiencia, sugiriendo el uso de máquinas "desorganizadas" que guardan una notable similitud con las redes neuronales modernas. Turing comprendió que la computación universal era el cimiento, pero sus estimaciones sobre el almacenamiento y la velocidad requerida reflejaban las restricciones tecnológicas de 1950

\section{La Revolución Computacional como Catalizador de las Redes Neuronales Profundas}
Para comprender la magnitud de la revolución actual, es imperativo remontarse a los primeros intentos de mecanizar el pensamiento. En 1956, durante la histórica conferencia de Dartmouth, los pioneros de la disciplina establecieron expectativas sumamente optimistas. Modelos tempranos, como el Perceptrón de Frank Rosenblatt a finales de la década de 1950, demostraron que una red artificial podía aprender a clasificar imágenes simples. No obstante, la euforia fue efímera

En 1969, Marvin Minsky y Seymour Papert publicaron una crítica devastadora que demostraba matemáticamente que los perceptrones de una sola capa eran incapaces de resolver funciones lógicas no lineales, como el operador XOR. Aunque la solución teórica a este problema era conocida —agregar múltiples capas ocultas (lo que hoy llamamos redes neuronales multicapa)—, Minsky y Papert señalaron, con razón, que las computadoras de la época carecían del poder de procesamiento necesario para manejar el inmenso volumen de cálculos que tales redes requerían.

\includepdf[pages=-]{lineaTiempo.pdf}
\newpage
Esta desconexión entre la viabilidad teórica y la inviabilidad computacional precipitó el primer gran "Invierno de la IA". Durante décadas, el financiamiento se secó y la investigación se volcó hacia los Sistemas Expertos, los cuales no requerían aprender de grandes volúmenes de datos, sino que operaban bajo reglas deductivas "si-entonces" (if-then) programadas manualmente. Estos sistemas funcionaban bien en entornos estructurados, pero fracasaban estrepitosamente al intentar emular el sentido común, la percepción visual o el lenguaje natural, tareas que inherentemente requieren análisis estadístico y probabilístico masivo

\subsection{La Eficiencia Algorítmica en la Era de la Escasez}

Durante los años de escasez computacional, el desarrollo de la IA dependió casi exclusivamente de la elegancia algorítmica. Antes de la llegada del hardware especializado, la optimización extrema era la única vía para ejecutar programas complejos. El diseño meticuloso de algoritmos requería un profundo entendimiento de la complejidad temporal y espacial (la notación Big O). Para exprimir cada ciclo de reloj del procesador, los ingenieros recurrían al uso intensivo de estructuras de datos avanzadas, como pilas, colas, conjuntos (sets) y multiconjuntos, así como a técnicas de programación dinámica implementadas en lenguajes de alto rendimiento y control estricto de memoria, como C o C++.

La programación de bajo nivel y el diseño algorítmico eficiente eran el único salvavidas frente a un hardware secuencial (CPU) de un solo núcleo. La Unidad Central de Procesamiento (CPU) es una maravilla de la ingeniería secuencial, optimizada para el cambio rápido de contexto y la lógica de control compleja. Sin embargo, el entrenamiento de una red neuronal multicapa no requiere decisiones lógicas altamente ramificadas; requiere la ejecución simultánea de millones de operaciones matemáticas simples (sumas y multiplicaciones de matrices). Las CPUs, a pesar de las optimizaciones algorítmicas mediante lenguajes compilados rigurosos, se convertían en un embudo irremediable para la propagación hacia atrás (backpropagation) requerida para ajustar los pesos en una red profunda.

\subsection{La Ley de Moore, el Fin del Escalado de Dennard y el ``Muro de la Memoria''}
Para comprender por qué las redes neuronales profundas tuvieron que esperar décadas para materializarse, es fundamental analizar las leyes físicas y económicas que gobernaron el desarrollo del silicio. Durante la segunda mitad del siglo XX, la industria de los semiconductores fue impulsada por dos principios rectores: la Ley de Moore, que predecía la duplicación del número de transistores en un microchip cada dos años, y el Escalado de Dennard (Dennard Scaling), formulado en 1974. El Escalado de Dennard establecía que, a medida que los transistores se hacían más pequeños, su densidad de potencia se mantenía constante, lo que permitía a los fabricantes aumentar la frecuencia de reloj (los gigahercios) sin que el chip se sobrecalentara de forma incontrolable.

Bajo este paradigma, los pioneros de la IA en los años 80 y 90 confiaban en que, eventualmente, las CPUs de un solo núcleo serían lo suficientemente rápidas para procesar redes neuronales complejas. Sin embargo, alrededor del año 2006, el Escalado de Dennard colapsó. Los transistores se volvieron tan diminutos que comenzaron a sufrir fugas de corriente a nivel cuántico, generando un calor masivo que impedía seguir aumentando la velocidad del reloj más allá de los 3.5 a 4.0 GHz. 

Al mismo tiempo, la industria chocó contra el "Muro de la Memoria" (Memory Wall). Aunque la velocidad de los procesadores había aumentado exponencialmente, la velocidad de la memoria RAM (la rapidez con la que se pueden enviar datos desde la memoria al procesador) crecía a un ritmo mucho más lento. En el contexto de una red neuronal, donde millones de "pesos" y "sesgos" deben ser actualizados constantemente durante el entrenamiento mediante la propagación hacia atrás, la CPU pasaba la mayor parte de su tiempo inactiva, esperando a que los datos llegaran desde la memoria. La inviabilidad de las redes neuronales en el pasado no solo se debía a una falta de "poder de cálculo" bruto, sino a una arquitectura física que era inherentemente hostil al movimiento masivo de matrices numéricas. Esta crisis obligó a la industria computacional a abandonar la carrera por la velocidad de un solo núcleo y abrazar el paralelismo multinúcleo, preparando, sin saberlo, el escenario perfecto para el surgimiento de las GPUs en el campo de la Inteligencia Artificial.

\subsection{El Despertar del Silicio: De la CPU al Paralelismo de la GPU}
El verdadero punto de inflexión no provino inicialmente de los laboratorios de IA, sino de la industria de los videojuegos. La necesidad de renderizar millones de píxeles en pantalla a sesenta fotogramas por segundo llevó al desarrollo de la Unidad de Procesamiento Gráfico (GPU). A diferencia de una CPU, que puede tener entre 4 y 16 núcleos muy potentes, una GPU moderna cuenta con miles de núcleos más simples diseñados para ejecutar la misma operación matemática simultáneamente a través de vastos conjuntos de datos (arquitectura SIMD: Single Instruction, Multiple Data).

A mediados de la década de 2000, los investigadores de IA se dieron cuenta de que las operaciones de álgebra lineal que fundamentan las redes neuronales profundas eran matemáticamente idénticas a las transformaciones matriciales utilizadas en los gráficos de computadora. El lanzamiento de plataformas de computación paralela permitió a los desarrolladores escribir código directamente para las GPUs. El impacto fue sísmico.

En 2012, una red neuronal convolucional profunda (CNN) llamada AlexNet aplastó a la competencia en el desafío de reconocimiento de imágenes ImageNet. Lo que antes habría tomado años entrenar en un clúster de CPUs, tomó solo unos días utilizando dos GPUs. La viabilidad computacional finalmente había alcanzado a la teoría de los años 80. Este avance marcó el fin definitivo del invierno de la IA y el inicio de la era del Deep Learning.

\subsection{La Revolución del Software Puente: GPGPU y el Ecosistema CUDA}
Si bien el hardware de las GPUs ofrecía el paralelismo masivo necesario para revivir a las redes neuronales, existía una barrera de abstracción monumental. A principios de la década del 2000, para que un investigador de Inteligencia Artificial pudiera utilizar una tarjeta gráfica para realizar cálculos matriciales, debía ''engañar'' al hardware. Tenía que traducir sus ecuaciones matemáticas de redes neuronales a un formato que la GPU entendiera: gráficos. Los tensores y matrices se codificaban como "texturas", y las operaciones matemáticas se escribían utilizando lenguajes de sombreado de píxeles (Pixel Shaders) como OpenGL o DirectX. Este proceso era un esfuerzo titánico, propenso a errores y limitaba drásticamente la experimentación teórica.

El verdadero catalizador que unió la viabilidad teórica con la capacidad computacional fue la introducción de la arquitectura CUDA (Compute Unified Device Architecture) por parte de NVIDIA en 2006. CUDA democratizó el hardware. Permitió a los ingenieros e investigadores escribir programas en C, C++ y posteriormente en Python, ejecutando algoritmos matemáticos genéricos directamente en los miles de núcleos de la GPU, sin necesidad de saber nada sobre la renderización de polígonos o píxeles. 

Este hito inauguró la era del GPGPU (General-Purpose computing on Graphics Processing Units). A partir de esta abstracción de bajo nivel, nacieron los frameworks de software de alto nivel que dominan el panorama actual, como TensorFlow (Google) y PyTorch (Meta/Facebook). Estas bibliotecas de software abstrajeron la complejidad del cálculo de gradientes y la asignación de memoria en la GPU. Sin esta capa de software puente, la capacidad bruta del silicio habría permanecido inaccesible para la mayoría de los matemáticos y científicos de datos. La viabilidad del Deep Learning no fue solo un triunfo del hardware, sino un triunfo de las APIs y los compiladores que supieron traducir la matemática pura en instrucciones máquina paralelas.

\subsection{La Arquitectura Transformer: Diseñada para la Escala Computacional}

El impacto del hardware en la arquitectura de software se hizo aún más evidente con el procesamiento del lenguaje natural. A principios de la década de 2010, los modelos dominantes eran las Redes Neuronales Recurrentes (RNN) y las redes de Memoria a Corto y Largo Plazo (LSTM). Aunque eficaces para secuencias, las RNN y LSTMs padecían de un defecto arquitectónico fundamental: procesaban los datos de forma estrictamente secuencial. Para predecir la palabra número cien de una oración, debían procesar las noventa y nueve anteriores una por una. Esto imposibilitaba el uso pleno del paralelismo masivo de las GPUs.

En 2017, la arquitectura Transformer alteró para siempre este paradigma. Basado en el mecanismo de "atención" (attention mechanism), el Transformer elimina la necesidad del procesamiento secuencial paso a paso. Permite que el modelo evalúe la relevancia de todas las palabras en una secuencia simultáneamente, calculando la relación entre cada token sin importar su distancia posicional

Desde una perspectiva de ingeniería, la mayor ventaja del Transformer no es solo su capacidad de capturar el contexto, sino su paralelizabilidad. Es una arquitectura diseñada específicamente para explotar la capacidad del hardware moderno. Esto ha permitido un escalado sin precedentes. Los modelos resultantes, como la serie GPT (Generative Pre-trained Transformer), han crecido exponencialmente en tamaño y capacidad. Mientras que GPT-1 (2018) contaba con 117 millones de parámetros, GPT-3 (2020) alcanzó los 175 mil millones, y GPT-4 superó la asombrosa cifra de 1.76 billones de parámetros

\subsection{Casos de Estudio Históricos: De LeNet-5 a los Modelos Fundacionales}
Para ilustrar concretamente la tesis de que la capacidad de cómputo actual habilitó arquitecturas previamente viables solo en papel, es esclarecedor contrastar dos hitos históricos: LeNet-5 y los modelos de visión contemporáneos. 

En 1998, el investigador Yann LeCun (galardonado posteriormente con el Premio Turing) desarrolló LeNet-5, una de las primeras arquitecturas funcionales de Redes Neuronales Convolucionales (CNN) diseñada para reconocer caracteres manuscritos, utilizada por bancos para leer cheques. Desde el punto de vista matemático, LeNet-5 posee exactamente las mismas capas fundamentales que las redes modernas de hoy: capas de convolución, capas de agrupación (pooling), funciones de activación no lineales (aunque usaba sigmoide en lugar de ReLU) y capas totalmente conectadas (fully connected). Teóricamente, el paradigma ya estaba completo. 

Sin embargo, LeNet-5 contenía apenas unos 60,000 parámetros. Entrenar este modelo diminuto en los procesadores de 1998 tomaba días. Si los investigadores de la época hubiesen intentado escalar esta misma red para clasificar imágenes a color de alta resolución en miles de categorías (como lo hace hoy la red ResNet o VGG), el tiempo de entrenamiento en el hardware de los años 90 habría tomado, literalmente, décadas e incluso siglos.

Hoy en día, un modelo fundacional de visión o lenguaje posee cientos de miles de millones de parámetros. La matemática de la propagación del error utilizando el descenso de gradiente estocástico (SGD) es idéntica a la que usó LeCun, y similar a la descrita teóricamente por Paul Werbos en 1980 con su Neocognitrón. Lo que cambió no fue la teoría subyacente de la neurobiología artificial, sino el hecho de que una GPU Nvidia H100 moderna puede realizar operaciones de precisión mixta a una velocidad de petaFLOPS (mil billones de operaciones por segundo). Este contraste histórico demuestra que las redes neuronales no estaban equivocadas en su concepción durante el siglo XX; simplemente estaban atrapadas en una época tecnológica prematura, esperando la llegada de la supercomputación distribuida.

\subsection{La Nube como Ecosistema Habilitador del Deep Learning}

Tener la arquitectura matemática y la idea de paralelización no es suficiente si no se cuenta con el espacio físico y lógico para albergarlo. Es crucial reflexionar sobre cómo la computación en la nube ha sido la columna vertebral silenciosa de esta revolución. El entrenamiento de un modelo de gran lenguaje (LLM) contemporáneo no ocurre en una supercomputadora monolítica, sino en vastos clústeres distribuidos compuestos por miles de aceleradores de hardware interconectados.

Para que la inteligencia artificial moderna pueda ser desarrollada, probada e implementada como un servicio consumible, se requiere una infraestructura de backend capaz de soportar cargas de trabajo masivas e impredecibles. Las capacidades de cómputo elástico, como el despliegue de múltiples instancias computacionales (por ejemplo, entornos EC2), brindan la fuerza bruta necesaria. Además, la implementación de grupos de autoescalado asegura que, ya sea en la fase de entrenamiento distribuido o durante los picos de inferencia generados por millones de usuarios, los recursos se adapten dinámicamente sin colapsar la infraestructura.

A la par del procesamiento, la gestión de datos es colosal. El entrenamiento de modelos como LLaMA o GPT requiere petabytes de texto, imágenes y código. Estos datos exigen soluciones de almacenamiento de objetos altamente durables y escalables (como arquitecturas S3). Además, cuando estas redes neuronales profundas trascienden la academia y se integran en aplicaciones interactivas —como agentes conversacionales o plataformas educativas—, la lógica del sistema, la gestión de perfiles de usuario, el historial de chats y el contexto de las sesiones dependen de sistemas de bases de datos relacionales altamente optimizados (utilizando motores robustos como PostgreSQL en servicios gestionados como RDS). La amalgama entre la fuerza bruta de las GPUs y la orquestación estructurada de la nube es lo que convierte a una red neuronal en un producto funcional y omnipresente en la sociedad actual.

\subsection{El Cuello de Botella de Von Neumann y la Memoria de Alto Ancho de Banda (HBM)}
Incluso con la computación en la nube y el paralelismo de las GPUs, la implementación de redes neuronales profundas gigantescas se topó con un desafío arquitectónico clásico: el Cuello de Botella de Von Neumann. En la arquitectura tradicional de computadoras propuesta por John von Neumann en 1945, la memoria y el procesador están físicamente separados. Para que el procesador realice un cálculo, primero debe buscar los datos en la memoria a través de un bus de comunicación.

En la era del Deep Learning, y particularmente con los Modelos de Lenguaje Extensos (LLMs) como GPT-4 o LLaMA, el cuello de botella más severo no es la velocidad de los núcleos de procesamiento, sino el ancho de banda de la memoria. Cuando se realiza la inferencia en un LLM (el proceso de generar texto palabra por palabra), el modelo entero debe cargarse repetidamente desde la memoria VRAM hacia los registros de cálculo. Si un modelo tiene 70 mil millones de parámetros, ocupa aproximadamente 140 Gigabytes de memoria en precisión de 16 bits. Mover 140 GB de datos físicamente a través de una placa de circuito en fracciones de milisegundo es un desafío físico y termodinámico extremo.

La solución tecnológica que realmente ha habilitado la viabilidad computacional de estos gigantescos modelos en los últimos cinco años ha sido la invención de la Memoria de Alto Ancho de Banda (HBM - High Bandwidth Memory). A diferencia de la memoria tradicional (GDDR), la HBM utiliza tecnología de apilamiento en 3D (chips de memoria colocados unos sobre otros) y se integra físicamente en el mismo encapsulado que el procesador de la GPU mediante un "interposer" de silicio. Esto reduce la distancia que viajan los datos de centímetros a milímetros, logrando anchos de banda de varios Terabytes por segundo. Sin esta innovación específica en la ciencia de los materiales y la arquitectura de memoria, los algoritmos de atención del Transformer serían inútiles, pues los núcleos de cálculo estarían perpetuamente inactivos esperando la llegada de los datos. Esta es otra prueba fehaciente de cómo la física del hardware dicta las fronteras de lo que la IA puede lograr teóricamente.


\subsection{El Costo Energético y el Límite Físico de la Computación}

Este avance monumental trae consigo desafíos tangibles que superan la barrera del software para adentrarse en la termodinámica. Las exigencias computacionales de los grandes modelos de lenguaje son titánicas. El entrenamiento de un modelo como GPT-3 requirió aproximadamente $3.14 \times 10^{23}$ operaciones de punto flotante por segundo (FLOPS), lo que se traduce en 355 años-GPU y un consumo de energía masivo que plantea serias preocupaciones ambientales respecto a la huella de carbono de la IA. El costo financiero de dicho entrenamiento se estimó en varios millones de dólares solo en tiempo de procesamiento.  

Este nivel de consumo está impulsando la siguiente fase de la revolución del hardware: los Circuitos Integrados de Aplicación Específica (ASIC), tales como las Unidades de Procesamiento Tensorial (TPU) de Google o los Procesadores Neuronales (NPU). A diferencia de las GPUs, que mantienen cierta versatilidad para el renderizado gráfico, los ASICs están esculpidos a nivel microscópico exclusivamente para acelerar las operaciones matriciales de las redes neuronales, maximizando el rendimiento por vatio.

Asimismo, estamos presenciando un retorno a la optimización algorítmica estricta. Así como en los años 80 y 90 se optimizaba el código en C++ por la falta de RAM y ciclos de CPU, hoy en día se investigan técnicas como la "cuantización" (reducir la precisión de los pesos de la red neuronal de 32 bits a 16, 8 o incluso 4 bits) y la destilación de conocimientos, con el fin de comprimir modelos mastodónticos para que puedan operar en hardware más modesto, o incluso en dispositivos móviles locales (Edge AI y TinyML).

\subsection{Consideraciones Filosóficas sobre el Paradigma Computacional}
Las herramientas analíticas concebidas por la teoría de la información y la computabilidad han permitido mapear un universo de modelos computacionales. La ciencia ha transitado de observaciones empíricas aleatorias a modelos no aleatorios que comprimen y explican la realidad de forma mecanicista. Como señalaba Turing, las máquinas operan en un espacio finito, pero su capacidad para imitar, adaptar y generar patrones complejos depende de su arquitectura y del entrenamiento que reciben.

Las redes neuronales profundas demuestran que el ''aprendizaje'' mecanizado requiere una escala de información y fuerza de procesamiento que emula, de cierta forma, la vasta densidad de interconexiones del cerebro biológico. Los primeros teóricos tenían los planos de la mente artificial, pero carecían de los ladrillos (datos) y los albañiles (ciclos de reloj). La capacidad de cómputo actual ha actuado como la materia prima para dar cuerpo a conceptos teóricos que llevaban décadas esperando su momento.


\subsection{El Horizonte Computacional: Computación Neuromórfica y Cuántica}
La actual viabilidad de las redes neuronales profundas ha llevado la física del silicio a sus límites fundamentales, impulsando la búsqueda de nuevos paradigmas computacionales que sostengan el futuro de la IA. Si la GPU y la memoria HBM desbloquearon la era actual, la próxima frontera busca resolver el colosal consumo energético mediante la Computación Neuromórfica y la Computación Cuántica.

La Computación Neuromórfica busca abandonar por completo la rígida arquitectura de Von Neumann y el procesamiento booleano digital, proponiendo chips que imitan físicamente la estructura del cerebro humano mediante "redes neuronales de impulsos" (Spiking Neural Networks - SNNs). En un procesador neuromórfico (como el chip Loihi de Intel o TrueNorth de IBM), no hay separación entre memoria y cálculo; las "sinapsis" de silicio retienen la información y procesan los picos analógicos de voltaje simultáneamente. Esto podría reducir el consumo energético de la IA generativa en órdenes de magnitud, pasando de requerir megavatios a consumir milivatios, imitando los apenas 20 vatios que consume un cerebro humano biológico para realizar razonamientos infinitamente más complejos.

Por otro lado, la Computación Cuántica promete resolver el problema dimensional del Machine Learning. Entrenar redes neuronales masivas es, en el fondo, encontrar el mínimo global en un espacio multidimensional de miles de millones de variables. Mientras el hardware clásico escala de forma lineal, los cúbits cuánticos escalan de forma exponencial gracias a la superposición y el entrelazamiento. Algoritmos emergentes de Quantum Machine Learning (QML) teorizan que los futuros ordenadores cuánticos podrán procesar el entrenamiento de redes neuronales profundas en tiempos logarítmicos, abriendo la puerta a modelos que hoy consideramos inalcanzables, como simulaciones completas de plegamiento de proteínas a nivel atómico. El ciclo histórico se repite: así como Turing y LeCun formularon teorías limitadas por el hardware de su tiempo, los teóricos actuales del QML están formulando algoritmos matemáticos que aguardan la viabilidad computacional del hardware cuántico comercial en las próximas décadas.

\subsection{El Efecto Lovelace y la Percepción de la Creatividad Mecánica}
La habilitación computacional de las redes neuronales profundas no solo ha resuelto problemas técnicos, sino que ha reabierto debates ontológicos fundamentales sobre la creatividad. Históricamente, la objeción de Lovelace planteaba que una máquina no puede originar nada por sí misma, sino que está limitada a realizar aquello para lo que ha sido programada\cite{natale2024}. Sin embargo, el advenimiento de modelos generativos impulsados por hardware masivo, como las Redes Generativas Antagónicas (GANs), ha transformado esta noción. 

El caso del sistema AICAN, entrenado con decenas de miles de imágenes de la historia del arte occidental, demuestra cómo la capacidad de cómputo actual permite a las redes neuronales aprender y generar obras que superan la mera replicación algorítmica\cite{natale2024}. Más allá de la ontología de la máquina, los académicos proponen hoy el ``Efecto Lovelace'', un concepto que desplaza el foco de lo que la máquina puede hacer hacia cómo los humanos perciben sus resultados\cite{natale2024}. Cuando una red neuronal como AICAN exhibe sus obras en galerías de arte contemporáneo, los espectadores experimentan el Efecto Lovelace al percibir genuina originalidad e intencionalidad en las piezas, a menudo siendo incapaces de distinguir si fueron creadas por un humano o por una inteligencia artificial\cite{natale2024}. Este fenómeno ilustra que el poder de cómputo no solo ha escalado la matemática, sino que ha alterado la relación semiótica entre humanos y máquinas, transformando modelos estadísticos en entidades percibidas como agentes creativos\cite{natale2024}.

\subsection{El Desafío de la Interpretabilidad y las Alucinaciones en Redes Profundas}
Paradójicamente, la misma explosión en la capacidad de cómputo que posibilitó el Deep Learning ha generado modelos tan masivos que resultan inescrutables. Las redes neuronales actuales, con miles de millones de parámetros, son frecuentemente criticadas por operar como ``cajas negras''\cite{zhang2023}. Aunque técnicas como los mapas de atención permiten visualizar en qué partes de una entrada se enfoca el modelo, estas herramientas no ofrecen una comprensión exhaustiva y determinista del proceso completo de toma de decisiones\cite{zhang2023}. 

Uno de los subproductos más críticos de esta masiva parametrización son las ``alucinaciones'' en la generación de lenguaje natural\cite{zhang2023}. A pesar de generar textos coherentes y gramaticalmente impecables, las redes neuronales pueden producir afirmaciones fácticas incorrectas debido al sobreajuste (overfitting) o a las limitaciones de sus datos de entrenamiento\cite{zhang2023}. La complejidad del modelo hace que la máquina pueda memorizar ruido estadístico en lugar de patrones reales\cite{zhang2023}. Además, debido a que estos modelos requieren cantidades colosales de datos de Internet, corren el riesgo de heredar, perpetuar e incluso amplificar sesgos raciales, de género y sociales presentes en la información original\cite{zhang2023, haenlein2019}. 

El reto actual, por lo tanto, trasciende la viabilidad computacional y se adentra en el terreno de la ética y la regulación. Distintas regiones están abordando este problema con enfoques divergentes; mientras Estados Unidos y China fomentan el desarrollo acelerado, la Unión Europea ha implementado regulaciones estrictas, como el GDPR, que imponen barreras sobre cómo se pueden procesar los datos, lo que plantea interrogantes sobre el equilibrio entre la innovación tecnológica y la privacidad individual\cite{haenlein2019}.

\subsection{Evolución del Paradigma de Turing: De la Imitación a la Colaboración}
El Test de Turing original de 1950 fue concebido en una época donde la informática era rudimentaria, sirviendo más como una bomba de intuición teórica (intuition pump) para la ingeniería conceptual del término ``inteligencia'' que como una métrica de rendimiento estricta\cite{milkowski2022}. El objetivo inicial era evadir la discusión filosófica sobre el pensamiento y reemplazarla con un desafío conductual operacional\cite{grosz2012}. Hoy en día, gracias a que el hardware permite el despliegue de agentes inteligentes complejos en la nube y en dispositivos móviles, la forma en que los humanos utilizan la computación ha cambiado drásticamente de un paradigma de uso individual a redes mixtas donde humanos y máquinas interactúan constantemente\cite{grosz2012}.

Ante esta realidad habilitada por la supercomputación, el desafío original de Turing resulta insuficiente. Como señala Barbara Grosz, las verdaderas capacidades cognitivas humanas están profundamente enraizadas en la cooperación, la intencionalidad y el trabajo en equipo, elementos que el Juego de la Imitación tradicional no evalúa\cite{grosz2012}. En la actualidad, el reto analítico no es crear un sistema que engañe a un juez haciéndose pasar por humano, sino desarrollar agentes de IA que puedan colaborar fluidamente a largo plazo en entornos dinámicos e inciertos como verdaderos compañeros de equipo\cite{grosz2012}. El hardware moderno ha permitido a la IA superar el rol de calculadora secuencial para convertirse en un participante activo capaz de formar modelos sobre las intenciones y planes de sus contrapartes biológicas\cite{grosz2012}.

\subsection{Complejidad Algorítmica y la Comprensión de la Naturaleza}
Finalmente, la materialización de las redes neuronales profundas gracias al cómputo masivo nos permite aproximarnos al universo desde la teoría de la información. El axioma de que ``la compresión es comprensión'' sostiene que encontrar modelos computacionales más pequeños y eficientes para explicar enormes cantidades de datos empíricos es, en esencia, la definición del avance científico\cite{zenil2021}. 

A través del aprendizaje profundo, el hardware moderno realiza exactamente este proceso: comprime petabytes de datos visuales, lingüísticos y físicos en redes de parámetros matemáticos. Las herramientas de probabilidad algorítmica y el Método del Teorema de Codificación (CTM) demuestran que gran parte del aparente ruido o aleatoriedad en el mundo natural oculta una estructura algorítmica subyacente\cite{zenil2021}. Las redes neuronales profundas, al ser universales y ahora computacionalmente viables, actúan como mecanismos masivos de descubrimiento causal; son capaces de desentrañar patrones y leyes no lineales en sistemas complejos que eran imposibles de procesar para la mente humana o para los sistemas de reglas rígidas del siglo pasado\cite{zenil2021}. Esto sugiere profundamente que la escalabilidad computacional actual no es solo una herramienta tecnológica, sino un lente cognitivo de magnitud sin precedentes para comprender la naturaleza algorítmica de nuestro propio universo\cite{zenil2021}.

\section{Conclusiones}
La afirmación de que la capacidad de cómputo actual ha habilitado teóricamente a las redes neuronales no es una subestimación; es la piedra angular del presente tecnológico. El viaje desde las cintas infinitas teóricas de la Máquina de Turing hasta los clústeres de servidores autoescalables que procesan billones de parámetros en paralelo ilustra una profunda lección de la ingeniería: la teoría marca el destino, pero es la capacidad de infraestructura y la innovación en el hardware la que construye el puente hacia él.

La superación de los inviernos de la IA se debió a un feliz encuentro histórico: la persistencia matemática de algoritmos como la propagación hacia atrás y la arquitectura Transformer, fusionada con el nacimiento del paralelismo masivo de las GPUs, el almacenamiento escalable en la nube y el uso de estructuras de bases de datos relacionales ultrarrápidas para organizar el contexto y los despliegues en producción. Las redes neuronales profundas siempre fueron computacionalmente válidas en el papel; fue nuestra especie la que necesitó casi medio siglo para fabricar el entorno físico y virtual capaz de soportar la inmensidad de su matemática. El desafío a futuro no radicará en si tenemos suficiente poder para sostener estos modelos, sino en cómo optimizar tanto el hardware especializado como el diseño algorítmico de bajo nivel para garantizar que este salto evolutivo sea energética y ecológicamente sostenible.


\vspace{0.5cm}
\begin{thebibliography}{99}

\bibitem{abramson2011} Abramson, D. (2011). \textit{Descartes' influence on Turing}. Studies in History and Philosophy of Science.
\bibitem{milkowski2022} Mi{\l}kowski, M. (2022). \textit{Turing's Conceptual Engineering}. Philosophies, 7(3), 69.
\bibitem{zhang2023} Zhang, E. Y., Cheok, A. D., Pan, Z., Cai, J., \& Yan, Y. (2023). \textit{From Turing to Transformers: A Comprehensive Review and Tutorial on the Evolution and Applications of Generative Transformer Models}. Sci, 5(4), 46.
\bibitem{zenil2021} Zenil, H. (2021). \textit{Compression is Comprehension: The Unreasonable Effectiveness of Digital Computation in the Natural World}. arXiv preprint arXiv:1904.10258.
\bibitem{haenlein2019} Haenlein, M., \& Kaplan, A. (2019). \textit{A Brief History of Artificial Intelligence: On the Past, Present, and Future of Artificial Intelligence}. California Management Review, 61(4), 5--14.
\bibitem{buchanan2005} Buchanan, B. G. (2005). \textit{A (Very) Brief History of Artificial Intelligence}. AI Magazine, 26(4), 53--60.
\bibitem{natale2024} Natale, S., \& Henrickson, L. (2024). \textit{The Lovelace Effect: Perceptions of Creativity in Machines}. New Media \& Society, 26(4), 1909--1926.
\bibitem{grosz2012} Grosz, B. J. (2012). \textit{What Question Would Turing Pose Today?} AI Magazine, 33(4), 73--81.

\end{thebibliography}

\end{document}